% =====================================================================
\section{Conclusion}

A reinforcement learning architecture that jointly revisits the yard
block and arrival time of external truck jobs in response to real-time
congestion has been developed. Proposal and acceptance are modelled as
independent policies and a central procedure applies the assignments in
one batch after checking mutual consent and resource constraints; the
uncertainty in the induced congestion effects is tackled with a
counterfactual cost signal from discrete-event simulations that differ
in a single action, while at operating time candidates are evaluated by
the trained networks alone. Three points follow. Spatial placement and
temporal adjustment are compared within one action range. Separating the
roles of proposing and accepting extends the logic of distributed
appointment negotiation to a real-time reallocation
policy~\cite{ref6}. Counterfactual comparison makes the
contribution of a single action learnable from a joint cost, evaluating
induced waiting and equipment flow in the same monetary unit as the
immediate movement cost; because the branch computation is separated
from operation, the decision path remains a small cost model and a
constraint check. The construction thereby connects arrival-time
negotiation~\cite{ref6}, yard crane operations~\cite{ref7,ref8},
and multi-agent counterfactual learning~\cite{ref10,ref11} into one
spatio-temporal assignment problem.

Observations in the literature-calibrated environment indicate that the
structure carries spatio-temporal reallocation through to execution, and
that arrival-time adjustment and learned selection produce a direction
of cost reduction under congested conditions, which supports the view
that real-time information at a smart port can lead to assignment
actions rather than monitoring alone~\cite{ref1}. The same direction
held when the action range was matched, although on the arrival-time
axis alone the rule was lower on more days and the proposed policy led
only in total.

\subsection{Limitations}
Six limitations bound these results. The environment is synthetic: its
scale and variation are calibrated to reported ranges, but no operating
record of a terminal was fitted, and the demand mix, the arrival-curve
coefficients, and two of the four cost coefficients are design values.
Congestion is rare on the stage---four of the 28 days---and that is
exactly where the effect lives, so the statistics rest on a small number
of expensive days; more congested days, rather than more days, is what
would settle the learning contribution. The contribution of the consent
structure itself has not been isolated: we confirmed that nothing is
committed on one-sided consent, but we did not measure what removing the
acceptance policy's veto would cost, so separating proposal from
acceptance remains a design choice supported by its operation rather
than by an ablation. Time-slot quotas were left open, so the
arrival-time results describe an upper range rather than an executable
schedule. The three-hour counterfactual horizon is derived from the
design---a 30-minute review window, up to 120 minutes of deferral, and a
30-minute arrival-pressure margin---rather than taken from a reported
value. Finally, the main comparison uses a single crane dispatching
rule, and because the ranking of those rules reverses with demand, the
size of the reallocation effect may shift under another one.

Each of these can be addressed one axis at a time: reflecting actual
appointment quotas would make the executable range of time adjustment
concrete, and repeating a wider variety of congested conditions would
delineate the contribution relative to the untrained model. Observation,
proposal, acceptance, commitment, execution, and counterfactual learning
have been organised into a single reproducible structure, which as
automated equipment and appointment information spread can serve as a
basis for connecting real-time congestion information to explainable
resource assignment decisions.
